\section{Implementation}

\label{sec:implementation}
%% ============================================================================

\subsection{Architecture Overview}

Lux-DAS is implemented as a module within the Lux node software (luxd), integrated
with the Platform VM for Primary Network blocks and available as a library for appchain VMs.

\begin{lstlisting}[language=Go, caption={Core DAS interfaces}]
// BlobEncoder encodes raw data into 2D RS extended matrix.
type BlobEncoder interface {
    Encode(data []byte) (*ExtendedMatrix, error)
    Decode(matrix *ExtendedMatrix) ([]byte, error)
}

// CommitmentScheme produces and verifies KZG commitments.
type CommitmentScheme interface {
    Commit(matrix *ExtendedMatrix) (*DACommitment, error)
    OpenCell(row, col int) (*CellProof, error)
    VerifyCell(commit *DACommitment, row, col int,
               value FieldElement, proof *CellProof) bool
}

// Sampler implements the DAS sampling protocol.
type Sampler interface {
    Sample(ctx context.Context, header *BlockHeader,
           securityParam int) (bool, error)
    GetCell(ctx context.Context, appchainID ids.ID,
            blobIdx uint32, row, col int) (*CellWithProof, error)
}
\end{lstlisting}

\subsection{Finite Field Arithmetic}

All field operations use the BLS12-381 scalar field
$\mathbb{F}_r$ with $r = 0x73eda753\ldots$ (a 255-bit prime). We employ the Montgomery
multiplication form for efficient modular arithmetic, with a specialized implementation for
the Number Theoretic Transform (NTT) used in polynomial evaluation and interpolation.

\begin{table}[H]
\centering
\caption{Field operation benchmarks (Apple M2, single core)}
\label{tab:field-ops}
\begin{tabular}{lrr}
\toprule
\textbf{Operation} & \textbf{Time (ns)} & \textbf{Throughput (Mops/s)} \\
\midrule
Field addition & 2.1 & 476 \\
Field multiplication (Montgomery) & 8.3 & 120 \\
Field inversion (Fermat) & 2{,}400 & 0.42 \\
NTT forward ($k=64$) & 48{,}000 & --- \\
NTT forward ($k=256$) & 310{,}000 & --- \\
\bottomrule
\end{tabular}
\end{table}

\subsection{KZG Trusted Setup}

Lux-DAS reuses the Ethereum KZG ceremony~\cite{ethereumceremony} structured reference string
(SRS), which supports polynomials of degree up to $4{,}096$. This is sufficient for blob
matrices up to $k = 2{,}048$ (supporting blobs up to $2048^2 \times 31 \approx 130$\,MiB).
Reusing an existing ceremony avoids the coordination overhead of a new trusted setup while
benefiting from the 141{,}416-participant contribution chain.

\subsection{Networking Layer}

Cell distribution and sampling queries use a dedicated gossip topic per appchain. The message
format is:

\begin{lstlisting}[language=Go, caption={Cell gossip message}]
type CellMessage struct {
    AppchainID  ids.ID          `json:"appchainId"`
    SlotNum   uint64          `json:"slot"`
    BlobIndex uint32          `json:"blobIndex"`
    Row       uint16          `json:"row"`
    Col       uint16          `json:"col"`
    Value     [32]byte        `json:"value"`
    RowProof  bls.G1Affine    `json:"rowProof"`
    ColProof  bls.G1Affine    `json:"colProof"`
}
\end{lstlisting}

Each cell message is 228 bytes (32-byte appchain ID + 8-byte slot + 4-byte index + 2+2 byte
coordinates + 32-byte value + 2 $\times$ 48-byte proofs + overhead). For a full blob with
$k=64$, disseminating all $4 \times 64^2 = 16{,}384$ cells requires approximately 3.7\,MiB
of gossip traffic, spread across all validators.

\subsection{Storage and Pruning}

Full nodes store the extended matrix $\hat{M}$ and all KZG proofs for a configurable retention
period (default: 30 days). After this period, only the DACommitment and the first $k$ rows of
the original data matrix are retained, enabling reconstruction but not sampling verification
for historical blocks. Light clients do not store any blob data; they rely on full nodes for
on-demand cell retrieval.

\subsection{Parallelism and Pipelining}

Encoding a blob is embarrassingly parallel across rows and columns. Our implementation uses
a worker pool of $\min(2k, \text{GOMAXPROCS})$ goroutines for the NTT evaluations, achieving
near-linear speedup up to 8 cores.

\begin{table}[H]
\centering
\caption{Encoding latency vs.\ core count ($k = 64$, blob $\approx 127$\,KiB)}
\label{tab:encoding-parallel}
\begin{tabular}{lrrrr}
\toprule
\textbf{Cores} & 1 & 2 & 4 & 8 \\
\midrule
Encoding (ms) & 48.2 & 25.1 & 13.4 & 7.8 \\
Speedup & 1.0$\times$ & 1.92$\times$ & 3.60$\times$ & 6.18$\times$ \\
\bottomrule
\end{tabular}
\end{table}


%% ============================================================================
